\documentclass{article}
\usepackage{amsmath,amsthm,amssymb}

\title{Type C Grothendieck}
\author{}
\date{}

\newtheorem{conjecture}{Conjecture}
\newtheorem{remark}{Remark}

\begin{document}
\maketitle

\section{Objects}

For a signed permutation \(\omega\), the type \(C\) stable Grothendieck
polynomial, or \(K\)-theoretic Stanley symmetric function of type \(C\), is
defined in the source note as
\[
  GC_\omega=\sum_{f\in F_\omega} \mathbf{x}^{\operatorname{wt}(f)}.
\]
Here \(F_\omega\) is the set of signed factorizations for \(\omega\).  Such a
factorization begins with a word in the ordered alphabet
\[
  \cdots < -3 < -2 < -1 < -0 < 0 < 1 < 2 < 3 < \cdots,
\]
where \(i\) and \(-i\), for \(i\ne 0\), represent the simple transposition
\(s_i\), while \(0\) and \(-0\) represent the type \(C\) generator \(s_0\).  The
word is subdivided into strictly increasing factors, and the weight records the
factor sizes.

The \(Q\)-Grothendieck function is
\[
  GQ_\lambda=\sum_{q\in Q_\lambda}\mathbf{x}^{\operatorname{wt}(q)},
\]
where \(Q_\lambda\) is the set of semistandard shifted set-valued tableaux of
\(Q\)-type.

\section{Conjectural expansion}

The source conjecture predicts a positive \(GQ\)-expansion
\[
  GC_\omega = \sum_{t\in T_\omega} GQ_{\operatorname{shape}(t)},
\]
where \(T_\omega\) is the set of conjectural type \(C\) unimodal Hecke tableaux
for \(\omega\).

The source note also states a stronger \(GR\)-model formulation
\[
  GC_\omega^+ = \sum_{t\in T_\omega} GR_{\operatorname{shape}(t)}.
\]
Here \(GC_\omega^+\) restricts the signed factorizations by requiring a unique
minimum-absolute-value entry in each factor, with positive sign.  The \(GR\)
model is the shifted set-valued tableau submodel in which the first \(i\) or
\(i'\), read left to right and bottom to top, is constrained to be \(i\).

\section{Three checked levels}

The source code tests three increasingly strong finite statements.

\begin{conjecture}[Basic count level]
For each signed permutation \(\omega\) and word length \(n\), Hecke words of
length \(n\) for \(\omega\) are counted by pairs \((t,r)\), where
\(t\in T_\omega\), \(r\) is a standard shifted set-valued tableau with \(n\)
entries, and \(t\) and \(r\) have the same shape.
\end{conjecture}

\begin{conjecture}[Strong count level]
The same count holds after restricting to Hecke words with no adjacent equal
letters and to standard shifted set-valued tableaux with no consecutive entries
in the same box.
\end{conjecture}

\begin{conjecture}[Peakset-preserving level]
The correspondence can also be made peakset-preserving: a peak
\(w_{i-1}<w_i>w_{i+1}\) in the Hecke word corresponds to a peak at \(i\) in the
recording shifted set-valued tableau.
\end{conjecture}

\begin{remark}
The source note explains that the peakset-preserving level would imply the
positive \(GQ\)-expansion for \(GC_\omega\), with coefficients counted by type
\(C\) unimodal tableaux of the corresponding shape.
\end{remark}

\section{Computer check}

The script \texttt{code/check\_type\_c\_grothendieck.py} is a curated combined
port of the three source scripts.  Its default run checks all three levels with
max word length \(4\) and generator indices \(0,\ldots,3\).  During curation,
all three levels were also checked with max word length \(5\) and generator
indices \(0,\ldots,3\).

The original source scripts use larger default parameters, including length
\(8\), largest generator \(4\) for the basic and strong checks, and length
\(11\), largest generator \(4\) for the strongest check.  The curated defaults
are smaller so that the item remains quick to run.

These computations are finite evidence and regression checks.  They do not
construct the conjectural bijections and they do not prove the conjectures.

\end{document}
